Elevation contracts/compresses the gastrointestinal linings and is key to lung pluging rates. 
Standard head-of-bed elevation for GERD relies primarily on gravity to reduce upward movement of gastric contents. 
An adjustable bed or wedge that elevates the upper body achieves this. 
Memory-foam or contoured surfaces (such as Tempur-Pedic) distribute pressure and conform to the body laying at/on compression and allowing your weight to compress against the barrier as your lungs actively plunge the depths of its chambersthe combination of:Gravity, and  
The body’s own weight creating compression against a supportive barrier

(example given: an elevated Tempur-Pedic or similar contoured/memory-foam bed).  According to the hypothesis, this gravity-plus-body-weight compression against the barrier is what is key to the lung-plunging rates used while mitigating the proposed dysrhythm.

Detailed statement of the thesis component regarding elevated plunging rates

The hypothesis reframes severe GERD as a primary neurological/neuromuscular disorder: a self-reinforcing rhythmic feedback loop involving vagal afferents from fundic stretch receptors, brainstem integration, vagal efferent-driven transient lower-esophageal-sphincter relaxations (TLESRs), and concurrent phrenic inhibition of the crural diaphragm. This coordinated barrier failure allows reflux; secondary sensitization then amplifies the loop.

Within this model, elevated plunging rates form a key mechanical-neurological intervention. The term refers to deliberate, full-depth respiratory excursions (“plunging” the lungs into the depths of their chambers) performed under specific postural conditions. These conditions are claimed to stabilize the rates and effectiveness of the plunging so that the respiratory dynamics directly counteract the phrenic inhibition and help interrupt the dysrhythmia.

Core elements of the elevated-plunging mechanism

Postural foundation (gravity + body-weight compression)  
   The upper body is elevated (adjustable bed, wedge, or equivalent incline). This is not simple flat lying. Gravity creates a passive downward vector that reduces retrograde fluid movement. Simultaneously, the patient’s own body weight presses into a contoured or memory-foam surface (example given: elevated Tempur-Pedic-type mattress). The conforming surface distributes pressure and is said to produce a compressive force against the gastrointestinal linings and the region of the esophagogastric junction / crural diaphragm.  

   According to the thesis, this dual force (gravity + guided body-weight compression) acts as a mechanical stabilizer or anchor for the thoracic cavity and the barrier structures.

Stabilization of lung-plunging rates  
   With the mechanical anchor in place, the patient performs active, deep lung plunging—full inspiratory and expiratory excursions that maximally expand and contract the lung chambers. The postural compression is claimed to:
   Limit excessive or dysregulated thoracic pressure swings,
   Maintain or restore appropriate crural-diaphragm tone during the respiratory cycle,
   Provide a consistent mechanical reference that allows the plunging rates to become regular and effective rather than chaotic or inhibited.

   In the model, the phrenic nerve’s normal rhythmic drive to the crural diaphragm is pathologically inhibited as part of the feedback loop. The stabilized plunging is proposed to re-engage or override that inhibition through a combination of mechanical support and altered afferent feedback (possibly involving lung-stretch receptors and the Hering-Breuer reflex or related vagal pathways).

Integration with the broader loop-interruption strategy  
   Elevated plunging is not presented in isolation. It is coordinated with temperature control (hot broth applied under the same postural conditions). The thermal stimulus is hypothesized to flood local receptors, alter smooth-muscle timing or afferent traffic, and act as a circuit-breaker. The mechanical-postural component supplies the stable platform that allows the respiratory dynamics to reinforce that interruption. Together they are said to reset the vagus–phrenic coordination failure.

Contrast with standard clinical view  
   Conventional guidance treats head-of-bed elevation purely as a passive gravitational method to keep gastric contents below the esophagogastric junction. The diaphragm is regarded mainly as a structural (and often compromised) component of the barrier.  

   The thesis expands this into a dynamic multi-force system:  
   Gravity reduces fluid reflux.  
   Body-weight compression against a contoured barrier stabilizes linings and crural support.  
   Stabilized lung-plunging rates actively modulate thoracic pressure and phrenic tone, thereby targeting the neurological dysrhythmia itself.

Self-test character  
   The intervention is framed as low-stakes and patient-observable, analogous to dietary elimination challenges. A person can adopt the elevated contoured posture, practice controlled deep plunging, optionally combine it with warm (not scalding) fluid, and monitor symptom change. Improvement is interpreted within the model as provisional support for the mechanical-neurological mechanism.

Summary of the elevated-plunging claim
Elevated plunging rates = deep, deliberate respiratory excursions performed while the upper body is elevated on a contoured surface so that gravity plus the patient’s own body weight create compressive stabilization of the gastrointestinal linings and crural-diaphragm region. This mechanical anchor is hypothesized to regularize thoracic pressure dynamics, preserve or restore crural tone against pathological phrenic inhibition, and thereby help interrupt the self-reinforcing vagus–phrenic feedback loop that the thesis identifies as the core neurological driver of severe GERD symptoms.

This component is an original synthesis within the holistic hypothesis. It draws on real anatomical relationships (vagal and phrenic innervation of the barrier, respiratory modulation of crural tone) while extending them into a specific, unvalidated causal and therapeutic sequence.

